\section{An angular estimate on a quaternion shell}\label{sec:angular}
Fix \(n,x,H\) as in Theorem~\ref{thm:main}. We first count coordinate
representations with an angular weight and prescribed local conditions.

\subsection{The norm shell}
Let \(R=\Z[\omega]\) and \(K=\Q(\omega)\). The Eisenstein lattice is
norm Euclidean, since every complex number is within squared
distance \(1/3\) of a lattice point. Write \(B=K+Kj\), where
\(j^2=-1\) and \(jz=\bar z j\), and put \(\Oq=R+Rj\).
Its reduced norm is \(\nrd(z+wj)=N(z)+N(w)\). The trace pairing on
\(1,\omega,j,\omega j\) has determinant \(-9\), so the order has
reduced discriminant \(3\). The algebra is definite and the number
of ramified places is even, so it is ramified at \(3\) and infinity.
Thus \(\Oq\) is maximal.

Rounding both complex coordinates gives a remainder of reduced norm
at most \(2/3\). In any nonzero integral right ideal, take a least
positive norm element \(\alpha\) and divide on the left by \(\alpha\).
The remainder in that ideal has norm smaller than \(\nrd(\alpha)\)
unless it vanishes. Every right ideal is therefore principal, as
are the fractional ideals by scaling. The unit group \(U\) consists
of the twelve elements of \(R^\times\cup R^\times j\). For
\(3\nmid m\), let \(\mathcal Q_m=\{\alpha\in\Oq:\nrd(\alpha)=m\}\).
At each \(\ell^e\Vert m\), integral right ideals correspond locally
to index-\(\ell^e\) sublattices of \(\Z_\ell^2\), of which there are
\(1+\ell+\cdots+\ell^e\). Each ideal has twelve generators of norm
\(m\), giving
\begin{equation}\label{eq:shell}
|\mathcal Q_m|=12\sigma(m),\qquad \sigma(m)=\sum_{d\mid m}d .
\end{equation}
We use this all-ideal convention throughout
\cite[Propositions~41.3.1 and~41.3.6]{Voight}.

The number of coordinates of a positive norm \(a\) is \(6f(a)\),
where \(f\) is multiplicative. Its local values are \(f(p^e)=e+1\)
at split primes \(p\equiv1\bmod3\), \(f(p^e)=1\) at \(3\), and
\(f(p^e)=\mathbf1_{2\mid e}\) at inert primes \(p\equiv2\bmod3\).
Indeed a split prime has two conjugate factors in \(R\), and a choice
of their exponents gives each possibility up to a unit. The ramified
prime in \(R\) above \(3\) is generated by \(1-\omega\) and has norm
\(3\). In particular, \(f(a)\le\tau(a)\), where \(\tau\) is the
divisor function. At an inert prime, \(R/pR\) is a field, so norm
divisibility by \(p\) is equivalent to coordinate divisibility by
\(p\). For even \(m\), both component norms have the same parity.
Both are even precisely on the shell scaled by \(2\) from \(m/4\).
The odd-component shell therefore has size
\begin{equation}\label{eq:odd-shell}
 R_o(m)=12\{\sigma(m)-\sigma(m/4)\},\qquad
 \frac{R_o(m)}m=18\prod_{\substack{p^e\Vert m\\p\ge5}}
                            \frac{\sigma(p^e)}{p^e},
\end{equation}
where \(\sigma(m/4)=0\) if \(4\nmid m\).

\subsection{Finite tests and the Hecke bound}
Use the finite group \(G\) obtained by multiplying the fixed factor
\(G_3=(\Oq/3\Oq)^\times/\F_3^\times\) and any factors
\(\PGL_2(\Z/p^a\Z)\) at \(p\ge5\), with \(a=1\) or \(2\).
The Hecke index in this subsection is coprime to the level.
Write \(\Gamma\) for the image of \(U\) in \(G\).
The kernel on \(U\) is \(\{\pm1\}\), as is seen at \(3\).
Each group has a norm squareclass map, using
determinant at split quaternion places. Both squareclasses occur.
At \(3\), for example, the norms of \(1\) and \(1+j\) give the two
classes. They are detected modulo \(p\) at every odd level prime.

For \(k\ge0\), let \(\rho_k=\Sym^{2k}\C^2\) denote the unitary
representation of \(\SU(2)\) on that space. For a real quaternion
of nonzero norm use \(u_\alpha=\alpha/\sqrt{\nrd(\alpha)}\in\SU(2)\),
under the usual identification of unit quaternions with
\(\SU(2)\). The action on order units descends to \(\Gamma\).
Use the space of functions \(F:G\to\Sym^{2k}\C^2\) satisfying
\(F(\gamma g)=\rho_k(\gamma)F(g)\), with squared norm summed once
per left coset. The all-ideal Hecke operator is
\begin{equation}\label{eq:hecke}
 (T_m F)(g)=\frac1{12}\sum_{\alpha\in\mathcal Q_m}
                 \rho_k(u_\alpha)^{-1} F(\alpha_{\rm fin} g),
\end{equation}
where \(\alpha_{\rm fin}\) is projective reduction in \(G\).
For clarity, class number one identifies this finite model with
the coefficient-valued automorphic space of \(PB^\times\) at the
indicated level \cite[Section~2, especially (4)--(6) and Proposition~2.5]{GV}.
At every finite place take the image of \(\Oq_p^\times\), replacing
it by the kernel of the indicated reduction at level primes.
Any finite adelic point can be moved into the product of the
unit groups by principalizing its ideal. The rational stabilizer
of that product is the image of \(U\). In fact a rational
representative of a stabilizer element is locally a scalar times
an order unit. One rational scalar cancels the valuations of all
the local scalars, so the representative becomes a unit everywhere. The right cosets
of the level subgroup then give the stated equivariance.

For \eqref{eq:hecke}, use one generator per right ideal, which
principalizes the finite component by \(\alpha^{-1}\) and gives the
archimedean multiplier \(\rho_k(u_\alpha)\). Changing the generator
by a right order unit cancels by left equivariance. Replacing each
\(\alpha\) by \(\bar\alpha\) gives \eqref{eq:hecke}; its denominator
is the number of generators per ideal. Conjugation also proves
self-adjointness. Schur orthogonality at infinity multiplies the
finite Hilbert norm by an overall constant in the automorphic
model, which cancels in operator norms.

We use the bound
\begin{equation}\label{eq:hecke-norm}
 \|T_m\|\le \tau(m)\sqrt m
\end{equation}
on the whole space for \(k>0\), and on the complement of the joint
norm-squareclass characters for \(k=0\). Here are the relevant
hypotheses for cuspidal transfer. A finite-dimensional smooth
irreducible local factor at a split place \(\ell\) is trivial on
the norm-one group, since its open normal kernel contains the
elementary unipotents by diagonal conjugation. Such a factor
cannot occur at positive weight. In fact \(\Oq[1/\ell]^1\) has
dense real image in \(\SU(2)\). Equation \eqref{eq:shell} gives
infinitely many norm-one elements \(\alpha/\ell^a\) of exact
denominator \(\ell^a\) by taking \(\nrd(\alpha)=\ell^{2a}\) and
omitting \(\ell\Oq\). There are at most \(12(2a+1)\) single-component
points of that norm, so these elements extend outside the Eisenstein
torus normalizer. A proper infinite closed subgroup of \(\SU(2)\)
normalizes its one-dimensional identity component. It contains the
order-three unit \(\omega\), which must then lie in that torus and
determines it as the Eisenstein torus. This contradicts the points
outside its normalizer. For a finite adelic representative \(g\),
impose the congruences \(g^{-1}\gamma g\) belonging to the level
away from \(\ell\). They cut out a finite-index subgroup of
\(\Oq[1/\ell]^1\), still dense by connectedness. At \(\ell\) the
discrepancy is absorbed by invariance under the local norm-one
group. Rational left invariance then forces the coefficient at
\(g\) to be fixed by that dense subgroup, which is impossible
in the nontrivial irreducible representation.

For \(k=0\), if the action at \(\ell\) is trivial on the norm-one
group, strong approximation for \(B^1\) away from
\(\{\infty,\ell\}\) makes it trivial on the whole finite adelic norm-one
group \cite[Main Theorem~28.5.3]{Voight}. Indeed, at a finite adelic
point \(g\), approximate \(ghg^{-1}\) off \(\ell\) for a norm-one
element \(h\) by a rational norm-one element, then use rational left
invariance and the local invariance at \(\ell\). Such functions
depend on the local norm squareclasses in \(G\), since a unit with
square norm can be rescaled to norm one. The joint sign characters
diagonalize this exceptional part: their eigenvalues are
\(\sigma(m)\) times the corresponding sign of \(m\), by
\eqref{eq:hecke}.

All other constituents have infinite-dimensional split local
factors. Jacquet--Langlands transfers them to cuspidal
\(\mathrm{GL}_2\) representations, of holomorphic weight \(2k+2\)
at infinity and unramified at primes of \(m\)
\cite[Theorem~14.4 and Proposition~5.20]{JL}. The unitary factor
in \eqref{eq:hecke} gives the classical Hecke eigenvalues divided
by \(m^k\). Equivalently the degree-\((\ell+1)\) spherical
operator at \(\ell\mid m\) has eigenvalue
\(\sqrt\ell(\alpha_\ell+\beta_\ell)\) in unitary Satake parameters.
Deligne gives \(|\alpha_\ell|=|\beta_\ell|=1\) and the trivial
central character gives \(\alpha_\ell\beta_\ell=1\)
\cite[Theorems~4.9 and~5.1]{DeligneModular}
\cite[Theorem~1.6]{DeligneWeil}. The all-ideal recurrence
\(T_{\ell^a}=T_\ell T_{\ell^{a-1}}-\ell T_{\ell^{a-2}}\) for \(a\ge2\)
then gives eigenvalue \(\ell^{a/2}\sum_{j=0}^a
\alpha_\ell^j\beta_\ell^{a-j}\), bounded by \((a+1)\ell^{a/2}\).
Multiplicativity proves \eqref{eq:hecke-norm} uniformly in the weight
and level.

\subsection{A finite test with a prescribed direction}
For a test \(b:G\to\C\) and a unit vector \(v\) in degree \(k\), define
\[
 F(g)=\frac1{|\Gamma|}\sum_{\gamma\in\Gamma}
                      \rho_k(\gamma)v\, b(\gamma^{-1}g).
\]
This is the orthogonal projection of \(b(g)v\) onto equivariant
functions in the full counting norm on \(G\). Its coset norm is at most
\(\|b\|_{\ell^2(G)}/\sqrt{|\Gamma|}\).
For another unit vector \(w\), reindexing the shell by the left
unit lift of each \(\gamma\) gives
\begin{equation}\label{eq:projection}
 12\langle w,(T_mF)(1)\rangle
 =\sum_{\alpha\in\mathcal Q_m}
       b(\alpha_{\rm fin})\,\langle w,\rho_k(u_\alpha)^{-1}v\rangle .
\end{equation}
Here the inner product is linear in its second variable.
The averaging factor cancels because for a lift \(\gamma_0\)
the substitution \(\alpha=\gamma_0\alpha'\) cancels
\(\rho_k(\gamma)\) against the archimedean multiplier.
This formula allows \(b\) to distinguish the two components.

Writing \(\alpha=z+yj\), put \(t_\alpha=(N(z)-N(y))/m\).
For \(v=w\) the unit vector proportional to \(X^kY^k\), the
matrix coefficient in \eqref{eq:projection} equals \(P_k(t_\alpha)\),
where \(P_k\) is the degree-\(k\) Legendre polynomial with \(P_k(1)=1\).
Indeed, the polynomial action gives the binomial formula
\[
 P_k(t_\alpha)=\sum_{j=0}^k(-1)^j\binom{k}{j}^2
               (N(z)/m)^{k-j}(N(y)/m)^j.
\]
Passing to \(u_\alpha^{-1}\) leaves the coefficient unchanged.

In degree zero, projecting onto the sign characters replaces \(b\)
by its average conditioned on the joint norm squareclasses of
\(m\). Indeed these characters form an orthogonal basis for
functions of the squareclasses, and the eigenvalue signs evaluate
them at \(m\). The averaging defining \(F\) preserves these means.
For a product of local tests the conditional mean is the product
of their conditional means. At each positive degree, evaluation
at one coset has norm at most one, so \eqref{eq:projection} is
bounded in absolute value by \(12\tau(m)\sqrt m\|b\|_{\ell^2(G)}\).
The same error bound applies to the scalar cuspidal part.

\subsection{The cost of a translated window}
Choose a fixed smooth \(0\le\phi\le1\) of compact support in
\((-1/2,1/2)\) and integral \(I_\phi>0\). Put \(h=H/n\),
\(b=M/n\), and \(\lambda=x/n\). We use the bump
\[
 W(t)=\phi\!\left(\frac{(1+t)/2-\lambda}{h}\right).
\]
As \(h\le b\), on its support \(1-t^2\asymp b\), and the
support has length at most \(2h\). The constant Legendre
coefficient of \(W\) is \(I_\phi h\).

\begin{lemma}\label{lem:bump}
The Legendre expansion \(W=\sum_{k\ge0}c_kP_k\) converges
uniformly and absolutely, with \(\sum_{k\ge1}|c_k|\ll_\phi\sqrt{b/h}\).
\end{lemma}
\begin{proof}
We use \(|P_k(t)|\ll k^{-1/2}(1-t^2)^{-1/4}\) for \(-1<t<1\).
For a short proof, let \(Y(u)=\sqrt{\sin u}\,P_k(\cos u)\) on
\(0<u<\pi\). The Legendre equation gives \(Y''+QY=0\) with
\(Q=(k+1/2)^2+1/(4\sin^2u)\).
The derivative of \(Y^2+(Y')^2/Q\) is
\(-Q'(Y')^2/Q^2\ge0\) up to \(\pi/2\). At \(\pi/2\) this energy
is \(O(k^{-1})\), by \(P_{2j}(0)=(-1)^j4^{-j}\binom{2j}{j}\),
\(P_{2j+1}(0)=0\) and \(P_k'(0)=kP_{k-1}(0)\).
Reflection gives the bound on both halves.

Write \(\mathcal L=-\frac d{dt}((1-t^2)\frac d{dt})\).
In the rescaled coordinate \(v=(t-2\lambda+1)/(2h)\), the operator
is \(b/h^2\) times a differential operator of order two with
coefficients and derivatives bounded on the support of \(\phi\).
Indeed its two coefficients in that scale are
\(-(1-t^2)/(4b)\) and \(th/b\). Integrating by parts twice
using \(\mathcal L P_k=k(k+1)P_k\), and also using the direct integral,
gives
\[
 |c_k|\ll_\phi h\sqrt{k}\,b^{-1/4}
       \min\!\left(1,\left(\frac{b}{h^2k^2}\right)^2\right).
\]
The factor \(\sqrt{k}\) includes the coefficient normalization
\((2k+1)/2\). Summing below and above \(K=\sqrt b/h\) gives
\(O_\phi(h b^{-1/4}K^{3/2})=O_\phi(\sqrt{b/h})\).
Since \(|P_k|\le1\) on \([-1,1]\), the series converges uniformly.
Its Legendre coefficients identify the sum with \(W\) by polynomial
density.
\end{proof}

Combining this lemma with \eqref{eq:projection}, a finite product
test of modulus at most one has weighted shell count
\begin{equation}\label{eq:angular}
 12\sigma(m)I_\phi h\prod_{p}\text{local conditional mean}
 +O_\phi\!\left(\sqrt{|G|}\,\tau(m)\sqrt m\sqrt{M/H}\right).
\end{equation}
The same \(W\) is used when the target is divided by a common factor.
Both component norms scale together and \(t_\alpha\) is unchanged.
